\subsection{Hille-Yosida 定理}\label{sec:hille-yosida-theorem}

\subsubsection{Hille-Yosida 定理}\label{subsec:hille-yosida-theorem}

前一节说明了每个强连续半群都有无穷小生成元, 也说明了如何利用半群求解初值问题. 从 PDE 的观点来看, 最重要的问题是如何识别那些作为 $C_0$ 半群无穷小生成元的算子. Hille-Yosida 定理回答这一问题.

\begin{Theorem}{Hille-Yosida}
\label{thm:hille-yosida-generation}
设 $A$ 为 Banach 空间 $X$ 中的算子. 则 $A$ 是满足 $\|T(t)\|\leq M\exp(\omega t)$ 的 $C_0$ 半群 $T(t)$ 的无穷小生成元, 当且仅当以下两个条件成立:\begin{enumerate}\item\label{hyp:hille-yosida-dense-closed} $\mathcal{D}(A)$ 稠密, 且 $A$ 为闭算子,\item\label{hyp:hille-yosida-resolvent-powers} 每个实数 $\lambda>\omega$ 都属于 $A$ 的预解集, 且对每个 $n\in\mathbb{N}$ 都有 $\|R_\lambda(A)^n\|\leq M/(\lambda-\omega)^n$.\end{enumerate}
\end{Theorem}

由下面的证明可见, 实际上每个满足 $\operatorname{Re}\lambda>\omega$ 的复数都属于预解集, 且 $\|R_\lambda(A)^n\|\leq M/(\operatorname{Re}\lambda-\omega)^n$.

\begingroup
\renewcommand{\endProof}{\par}
\begin{Proof}
已经知道第一个条件是必要的. 为看出第二个条件的必要性, 回顾对任意 $C_0$ 半群, 都存在常数 $M$ 和 $\omega$, 使 $\|T(t)\|\leq M\exp(\omega t)$. 现在考虑积分
\begin{equation}\label{eq:hille-yosida-resolvent-power-laplace-integral}
I_n(\lambda)x=\dfrac{1}{(n-1)!}\int_0^\infty t^{n-1}e^{-\lambda t}T(t)x\,\mathrm{d}t.
\end{equation}
显然, 当 $\operatorname{Re}\lambda>\omega$ 时, $I_n(\lambda)$ 有良好定义, 且
\begin{equation}\label{eq:hille-yosida-laplace-integral-norm-bound}
\|I_n(\lambda)\|\leq\dfrac{1}{(n-1)!}\int_0^\infty t^{n-1}e^{-\operatorname{Re}\lambda t}Me^{\omega t}\,\mathrm{d}t=\dfrac{M}{(\operatorname{Re}\lambda-\omega)^n}.
\end{equation}
若 $x\in\mathcal{D}(A)$, 则对 $n>1$, 有\[\begin{aligned}
I_n(\lambda)Ax&=\dfrac{1}{(n-1)!}\int_0^\infty t^{n-1}e^{-\lambda t}T(t)Ax\,\mathrm{d}t\\
&=\dfrac{1}{(n-1)!}\int_0^\infty t^{n-1}e^{-\lambda t}\dfrac{\mathrm{d}}{\mathrm{d}t}(T(t)x)\,\mathrm{d}t\\
&=-\dfrac{1}{(n-2)!}\int_0^\infty t^{n-2}e^{-\lambda t}T(t)x\,\mathrm{d}x\\
&\quad+\dfrac{\lambda}{(n-1)!}\int_0^\infty t^{n-1}e^{-\lambda t}T(t)x\,\mathrm{d}x\\
&=-I_{n-1}(\lambda)x+\lambda I_n(\lambda)x.
\end{aligned}\]当 $n=1$ 时, 则得到 $I_1(\lambda)Ax=-x+\lambda I_1(\lambda)x$. 也可将 $(T(h)-I)/h$ 作用于 \eqref{eq:hille-yosida-resolvent-power-laplace-integral}, 并进行类似的涉及 ``分部差分\textquotesingle\textquotesingle\ 的计算. 在极限 $h\to0$ 下得到\[\begin{aligned}
AI_n(\lambda)x&=-I_{n-1}(\lambda)x+\lambda I_m(\lambda)x,\ x\in X,\ n>1,\\
AI_1(\lambda)x&=-x+\lambda I_1(\lambda)x,\ x\in X.
\end{aligned}\]总之, 对 $n>1$, 有 $I_n(\lambda)(A-\lambda I)=(A-\lambda I)I_n(\lambda)=-I_{n-1}(\lambda)$, 且 $I_1(\lambda)(A-\lambda I)=(A-\lambda I)I_1(\lambda)=-I$. 这显然蕴含 $R_\lambda(A)$ 存在, 且 $I_n(\lambda)=(-R_\lambda(A))^n$.
\end{Proof}
\endgroup

现在转向充分性的证明. 为此考虑算子
\begin{equation}\label{eq:hille-yosida-implicit-resolvent-power-approximant}
U_n(t)=\left(I-\dfrac{t}{n}A\right)^{-n}.
\end{equation}
由条件~\ref{hyp:hille-yosida-resolvent-powers} 可知, 当 $n/t>\omega$ 时 $U_n(t)$ 有良好定义, 特别地, 对足够大的 $n$ 如此. 此外, 当 $n\to\infty$ 时, 算子 $U_n(t)$ 一致有界. 对 $t>0$, 关于 $t$ 微分得到
\begin{equation}\label{eq:hille-yosida-approximant-time-derivative}
\dot U_n(t)=A\left(I-\dfrac{t}{n}A\right)^{-n-1}
\end{equation}
这里取算子范数拓扑的意义. 断言当 $t\to0$ 时 $U_n(t)\to I$ 强收敛. 为此, 只需证明 $(I-\frac{t}{n}A)^{-1}$ 强收敛到 $I$; 见 Problem~\ref{ex:strong-operator-products-converge}. 由于已知 $(I-\frac{t}{n}A)^{-1}$ 在 $t\to0$ 时一致有界, 只需对某稠密集中的 $u$ 证明 $(I-\frac{t}{n}A)^{-1}u$ 收敛. 现在注意到, 对 $u\in\mathcal{D}(A)$,
\begin{equation}\label{eq:hille-yosida-resolvent-initial-strong-limit-bound}
\left\|\left(I-\dfrac{t}{n}A\right)^{-1}u-u\right\|=\left\|\dfrac{t}{n}\left(I-\dfrac{t}{n}A\right)^{-1}Au\right\|\leq Ct\|Au\|,
\end{equation}
其在 $t\to0$ 时趋于零.

断言当 $n\to\infty$ 时, $U_n(t)$ 有强极限. 为此, 对稠密集 $\mathcal{D}(A^2)=(A-\lambda)^{-1}\mathcal{D}(A)$ 中的 $u$ 估计 $U_n(t)u-U_m(t)u$, 其中 $\operatorname{Re}\lambda>\omega$. 注意到\[\begin{aligned}
U_n(t)u-U_m(t)u&=\lim_{\epsilon\to0}\int_\epsilon^{t-\epsilon}\dfrac{\mathrm{d}}{\mathrm{d}s}[U_m(t-s)U_n(s)u]\,\mathrm{d}s\\
&=\lim_{\epsilon\to0}\int_\epsilon^{t-\epsilon}[-\dot U_m(t-s)U_n(s)u+U_m(t-s)\dot U_n(s)u]\,\mathrm{d}s.
\end{aligned}\]利用 \eqref{eq:hille-yosida-approximant-time-derivative}, 经过一些代数运算得到
\begin{equation}\label{eq:hille-yosida-two-approximants-difference-integral}
\begin{aligned}U_n(t)u-U_m(t)u=\int_0^t&\left(\dfrac{s}{n}-\dfrac{t-s}{m}\right)\left(I-\dfrac{t-s}{m}A\right)^{-m-1}\\&\left(I-\dfrac{s}{n}A\right)^{-n-1}A^2u\,\mathrm{d}s.\end{aligned}
\end{equation}
于是
\begin{equation}\label{eq:hille-yosida-approximants-cauchy-bound}
\begin{aligned}\|U_n(t)u-U_m(t)u\|&\leq C\|A^2u\|\int_0^t\dfrac{s}{n}+\dfrac{t-s}{m}\,\mathrm{d}s\\&=\dfrac{Ct^2}{2}\left(\dfrac{1}{n}+\dfrac{1}{m}\right)\|A^2u\|.\end{aligned}
\end{equation}
因此 $U_n(t)$ 在 $n\to\infty$ 时的强极限存在, 将其记为 $T(t)$. 由于
\begin{equation}\label{eq:hille-yosida-approximant-exponential-growth-bound}
\|U_n(t)\|\leq M\left(1-\dfrac{t}{n}\omega\right)^{-n},
\end{equation}
在极限 $n\to\infty$ 下得到 $\|T(t)\|\leq M\exp(\omega t)$.

还需证明 $T(t)$ 确为半群, 且 $A$ 为其无穷小生成元. 由 \eqref{eq:hille-yosida-approximants-cauchy-bound} 容易看出, $U_n(t)u$ 到 $T(t)u$ 的收敛对有限区间中的 $t$ 确实一致. 由于 $U_n(t)u$ 对每个 $n$ 都连续, $T(t)u$ 对每个 $u\in\mathcal{D}(A^2)$ 都连续. 由于 $\mathcal{D}(A^2)$ 稠密且 $\|T(t)\|$ 在有限区间上有界, 可知 $T(t)u$ 对每个 $u\in X$ 都连续. 因而 $T(t)$ 强连续. 又由于对每个 $n$ 都有 $U_n(0)=I$, 显然有 $T(0)=I$. 此外, 与导出 \eqref{eq:hille-yosida-two-approximants-difference-integral} 相同的计算给出
\begin{equation}\label{eq:hille-yosida-semigroup-composition-derivative-bound}
\begin{aligned}\dfrac{\mathrm{d}}{\mathrm{d}s}[U_n(t-s)U_n(s)u]&=\dfrac{2s-t}{n}\left(I-\dfrac{t-s}{n}A\right)^{-n-1}\\&\quad\left(I-\dfrac{s}{n}A\right)^{-n-1}A^2u\end{aligned}
\end{equation}
对 $u\in\mathcal{D}(A^2)$ 成立. 令 $n\to\infty$, 可知对 $u\in\mathcal{D}(A^2)$, $T(t-s)T(s)$ 与 $s$ 无关; 再由稠密性, 对 $u\in X$ 也如此. 这就是半群性质.

\begingroup
\renewenvironment{Proof}{\par}{\par\nobreak\hfill$\blacksquare$\par\vspace{0.1cm}}
\begin{Proof}
在 \eqref{eq:hille-yosida-approximant-time-derivative} 中令 $n\to\infty$, 得到
\begin{equation}\label{eq:hille-yosida-limit-generator-initial-derivative}
\left.\dfrac{\mathrm{d}}{\mathrm{d}t}T(t)u\right|_{t=0}=Au
\end{equation}
对每个 $u\in\mathcal{D}(A)$ 成立; 因而该半群的无穷小生成元是 $A$ 的延拓. 这一延拓不可能是真延拓, 因为根据 $\|T(t)\|$ 的增长估计以及本证明的必要性部分, 该无穷小生成元在 $\lambda>\omega$ 时有预解算子. 这一预解算子必须与 $A$ 的预解算子一致, 因为定义在整个空间上的有界算子不能有真延拓.
\end{Proof}
\endgroup

上述证明不仅给出半群的存在性, 还给出逼近半群的实用方法. 有
\begin{equation}\label{eq:semigroup-implicit-resolvent-approximation-formula}
T(t)=\lim_{n\to\infty}\left(I-\dfrac{t}{n}A\right)^{-n}.
\end{equation}
这对应于求解方程 $\dot u=Au$ 的差分格式
\begin{equation}\label{eq:semigroup-implicit-euler-difference-scheme}
\dfrac{u(t+h)-u(t)}{h}=Au(t+h)
\end{equation}
即所谓隐式 Euler 格式. 它也对应于矩阵指数的公式 \eqref{eq:matrix-exponential-resolvent-power-limit}. 现在完全有理由不再对半群及其无穷小生成元采用分开的记号.

\begin{Definition}
\label{def:semigroup-exponential-notation}
设 $A$ 为 $X$ 上有界线性算子 $C_0$ 半群的无穷小生成元. 则用 $\exp(At)$, $t\geq0$, 表示由 $A$ 生成的半群.
\end{Definition}

还引入以下方便的记号.

\begin{Definition}
\label{def:semigroup-generator-class}
若 $X$ 中线性算子 $A$ 满足 Theorem~\ref{thm:hille-yosida-generation} 的假设, 则称其属于 $G(M,\omega)$.
\end{Definition}

\subsubsection{Lumer-Phillips 定理}\label{subsec:lumer-phillips-theorem}

Hille-Yosida 定理的条件在应用中仍不易验证, 因为它们要求对预解算子的所有幂给出界. 不过, 当 $M=1$ 时情形简单得多: 若 $\|R_\lambda(A)\|\leq(\lambda-\omega)^{-1}$, 则显然 $\|R_\lambda(A)^n\|\leq(\lambda-\omega)^{-n}$. 鉴于这一情形的重要性, 作如下定义.

\begin{Definition}
\label{def:quasicontraction-contraction-semigroup}
若对某个 $\omega$ 有 $\|T(t)\|\leq\exp(\omega t)$, 则称 $C_0$ 半群 $T(t)$ 为\textbf{拟压缩半群}. 若 $\|T(t)\|\leq1$, 则称它为\textbf{压缩半群}.
\end{Definition}

显然, 若 $A$ 生成拟压缩半群, 则 $A-\omega I$ 生成压缩半群. 实际上, 每个半群都可以变为拟压缩半群.

\begin{Theorem}
\label{thm:semigroup-equivalent-quasicontraction-norm}
设 $T(t)$ 为 $X$ 上满足
\begin{equation}\label{eq:semigroup-equivalent-quasicontraction-norm-growth-hypothesis}
\|T(t)\|\leq M\exp(\omega t).
\end{equation}
的 $C_0$ 半群. 则 $X$ 上存在一个等价范数, 使在与 $X$ 上该新范数对应的算子范数下, 有 $\|T(t)\|\leq\exp(\omega t)$.
\end{Theorem}

\PDEIndexMark{idx:lumer-phillips-criterion-introduction}略去该定理的证明; 例如可见 \cite{Pa}. 在许多应用中, 寻找适当的等价范数比直接验证 Hille-Yosida 定理的假设更有利. 以下称为 Lumer-Phillips 定理的结果, 给出拟压缩半群生成元的一个实用判据.

\begin{Theorem}{Lumer-Phillips}
\label{thm:lumer-phillips-generation}
设 $H$ 为 Hilbert 空间, $A$ 为 $H$ 中满足以下条件的线性算子:\begin{enumerate}\item\label{hyp:lumer-phillips-dense-domain} $\mathcal{D}(A)$ 稠密.\item\label{hyp:lumer-phillips-quasidissipativity} 对每个 $x\in\mathcal{D}(A)$, 都有 $\operatorname{Re}(x,Ax)\leq\omega(x,x)$.\item\label{hyp:lumer-phillips-surjectivity} 存在 $\lambda_0>\omega$, 使 $A-\lambda_0I$ 为满射.\end{enumerate}则 $A$ 生成拟压缩半群, 且 $\|\exp(At)\|\leq\exp(\omega t)$.
\end{Theorem}

\begingroup
\renewcommand{\endProof}{\par\nobreak\hfill$\blacksquare$\par\vspace{0.1cm}}
\begin{Proof}
对 $\lambda>\omega$, 有
\begin{equation}\label{eq:lumer-phillips-resolvent-lower-bound}
\|(A-\lambda I)x\|\|x\|\geq\operatorname{Re}(x,(\lambda I-A)x)\geq(\lambda-\omega)(x,x),
\end{equation}
即有 $\|(A-\lambda I)x\|\geq(\lambda-\omega)\|x\|$. 若能证明 $A-\lambda I$ 对每个 $\lambda>\omega$ 都是满射, 则这就蕴含 $R_\lambda(A)$ 存在, 且 $\|R_\lambda(A)\|\leq(\lambda-\omega)^{-1}$. 特别地, 有 $\lambda_0\in\rho(A)$, 从而 $A$ 为闭算子. 对 $\lambda>\omega$, $A-\lambda I$ 有有界逆, 因而其值域必为闭集. 但这意味着 $A-\lambda I$ 为半 Fredholm 算子, 且其指标必在 $(\omega,\infty)$ 上为常数.
\end{Proof}
\endgroup

\begin{Example}
\label{ex:self-skew-adjoint-semigroup-generation}
任何谱有上界的自伴算子都生成拟压缩半群. 任何斜自伴算子都生成压缩半群.
\end{Example}

\begin{Example}
\label{ex:absorbing-translation-dissipative-generator}
设 $H=L^2(0,1)$, 令 $Au=u^{\prime}$, 其定义域为 $\mathcal{D}(A)=\{u\in H^1(0,1)\mid u(1)=0\}$. 显然 $\mathcal{D}(A)$ 稠密, 且容易证明 $A$ 的谱为空. 此外, 有
\begin{equation}\label{eq:absorbing-translation-generator-dissipation-identity}
(u,Au)=\int_0^1u(x)u^{\prime}(x)\,\mathrm{d}x=-\dfrac{1}{2}(u(0))^2\leq0
\end{equation}
对每个 $u\in\mathcal{D}(A)$ 成立; 因而 $A$ 生成压缩半群. 实际上, $A$ 生成的半群就是 Example~\ref{ex:absorbing-translation-minus-infinite-type} 中给出的半群; 参见 Problem~\ref{ex:absorbing-translation-generator-characterization}.
\end{Example}

\begin{Definition}
\label{def:quasidissipative-dissipative-quasi-m}
Hilbert 空间中满足 Theorem~\ref{thm:lumer-phillips-generation} 的条件~\ref{hyp:lumer-phillips-quasidissipativity} 的算子称为\textbf{拟耗散的}; 若 $\omega=0$, 则称其为\textbf{耗散的}. 满足条件~\ref{hyp:lumer-phillips-quasidissipativity} 和~\ref{hyp:lumer-phillips-surjectivity} 的算子称为\textbf{拟 m 耗散的}.
\end{Definition}

\Needspace{5\baselineskip}\textbf{Problems}

\begin{enumerate}[label=\thesection.\arabic*.,ref=\thesection.\arabic*,start=9,series=book-problems]

\item\label{ex:hille-yosida-heat-dirichlet-existence} 证明热方程\[u_t=-Lu+f(t),\quad u(0)=u_0,\]的解的存在性, 其中 $L$ 为 \eqref{eq:second-order-garding-divergence-operator} 定义的带 Dirichlet 边界条件的二阶椭圆算子. 指明数据和解应属于的空间.

\item\label{ex:strong-operator-products-converge} 设 $A_k\in\mathcal{L}(X,Y)$, $B_k\in\mathcal{L}(Y,Z)$, 并假设当 $k\to\infty$ 时, $A_k\to A$, $B_k\to B$ 强收敛. 证明 $B_kA_k\to BA$ 强收敛.

\item\label{ex:quasicontraction-bounded-operator-perturbation} 设 $A\in G(1,\omega)$, 且 $B$ 有界. 证明 $A+B\in G(1,\omega+\|B\|)$.

\item\label{ex:translation-opposite-boundary-not-generator} 在 Example~\ref{ex:absorbing-translation-dissipative-generator} 中, 将边界条件 $u(1)=0$ 替换为条件 $u(0)=0$. 证明所得算子\emph{不}生成 $C_0$ 半群.

\item\label{ex:opposite-generators-inverse-semigroup-group} 若 $A$ 和 $-A$ 都满足 Hille-Yosida 定理的假设, 证明 $\exp(At)$ 和 $\exp((-A)t)$ 互为逆. 此时自然可将 $\exp((-A)t)$ 写成 $\exp(-At)$. 证明算子 $\exp(At)$, $t\in\mathbb{R}$ 构成一个群.

\item\label{ex:skew-adjoint-semigroup-unitary} 设 $A$ 为斜自伴算子. 证明 $\exp(At)$ 是酉的.

\end{enumerate}

\subsection{对 PDE 的应用}\label{sec:semigroup-pde-applications}

\subsubsection{对称双曲系统}\label{subsec:semigroup-symmetric-hyperbolic-systems}

设 $\bm{A}^i(\bm{x})$ 为对 $\bm{x}\in\mathbb{R}^m$ 定义的对称 $p\times p$ 矩阵. 在函数空间 $L^2(\mathbb{R}^m)$ 中考虑算子
\begin{equation}\label{eq:semigroup-symmetric-hyperbolic-spatial-operator}
\bm{A}\bm{u}(\bm{x})=\bm{A}^i(\bm{x})\dfrac{\partial\bm{u}}{\partial x_i}
\end{equation}
假设 $\bm{A}^i\in C_b^1(\mathbb{R}^m)$. 刻画 $\bm{A}$ 的定义域并非易事. 因而采取一种回避的办法: 先在 $H^1(\mathbb{R}^m)$ 上定义 $\bm{A}$, 再取闭包. 断言

\begin{Theorem}
\label{thm:symmetric-hyperbolic-semigroup-generation}
$\bm{A}$ 为某个 $C_0$ 半群的无穷小生成元.
\end{Theorem}

\begingroup
\renewcommand{\endProof}{\par}
\begin{Proof}
将使用 Lumer-Phillips 定理. 由于 $\bm{A}$ 的定义域包含 $H^1(\mathbb{R}^m)$, 显然 $\bm{A}$ 稠密定义. 此外, 对 $\bm{u}\in H^1$, 分部积分得到\[\begin{aligned}
(\bm{u},\bm{A}\bm{u})&=\int_{\mathbb{R}^m}\bm{u}(\bm{x})\cdot\bm{A}^i(\bm{x})\dfrac{\partial\bm{u}}{\partial x_i}\,\mathrm{d}\bm{x}\\
&=-\dfrac{1}{2}\int_{\mathbb{R}^m}\bm{u}\cdot\dfrac{\partial\bm{A}^i}{\partial x_i}\bm{u}\,\mathrm{d}\bm{x}.
\end{aligned}\]因此, Lumer-Phillips 定理的第二个条件对每个 $\bm{u}\in H^1(\mathbb{R}^m)$ 都成立; 通过取极限可证明, 它对每个 $\bm{u}\in\mathcal{D}(\bm{A})$ 也成立.
\end{Proof}
\endgroup

主要工作是验证第三个条件, 即对足够大的 $\lambda\in\mathbb{R}$, $\bm{A}-\lambda I$ 为满射. 由 Lumer-Phillips 定理的证明可见, 只需证明 $\bm{A}-\lambda I$ 的值域稠密. 为此, 先考虑一个新算子 $\widetilde{\bm{A}}$, 除定义域外它与 $\bm{A}$ 相同. 具体地, 取 $\widetilde{\bm{A}}$ 的定义域为所有满足如下条件的 $\bm{u}\in L^2$ 构成的集合: $\bm{A}^i\partial\bm{u}/\partial x_i$ (按 $H^{-1}(\mathbb{R}^m)$ 中元素理解) 也属于 $L^2(\mathbb{R}^m)$. 显然 $\widetilde{\bm{A}}$ 为 $\bm{A}$ 的延拓. 再用 Galerkin 论证证明, 对足够大的 $\lambda$, $\widetilde{\bm{A}}-\lambda I$ 为满射. 为此, 设 $\{\bm{\phi}_n\}$, $n\in\mathbb{N}$, 为 $H^1(\mathbb{R}^m)$ 中的完备正交系. 设 $X_n$ 为 $\{\bm{\phi}_1,\ldots,\bm{\phi}_n\}$ 的张成, $P_n$ 为从 $L^2(\mathbb{R}^m)$ 到 $X_n$ 的正交投影. 现在考虑方程
\begin{equation}\label{eq:symmetric-hyperbolic-resolvent-galerkin-equation}
(\bm{\phi},(\bm{A}-\lambda I)\bm{u}_n)=(\bm{\phi},\bm{f})\quad\forall\bm{\phi}\in X_n,
\end{equation}
其中寻求解 $\bm{u}_n\in X_n$. 利用 $\bm{A}$ 拟耗散这一事实, 容易看出这样的解 $\bm{u}_n$ 存在, 且 $\bm{u}_n$ 的 $L^2$ 范数有不依赖于 $n$ 的界. 因而 $\bm{u}_n$ 的某个子序列弱收敛, 且容易证明其极限为 $(\widetilde{\bm{A}}-\lambda I)\bm{u}=\bm{f}$ 的解.\par 还需证明 $\widetilde{\bm{A}}$ 确实就是 $\bm{A}$. 即给定满足 $\widetilde{\bm{A}}\bm{u}=\bm{f}$ 的 $\bm{u},\bm{f}\in L^2(\mathbb{R}^m)$, 需要找到 $\bm{u}_n\in H^1(\mathbb{R}^m)$, $\bm{f}_n\in L^2(\mathbb{R}^m)$, 使 $\bm{A}\bm{u}_n=\bm{f}_n$, 且 $\bm{u}_n\to\bm{u}$, $\bm{f}_n\to\bm{f}$ 在 $L^2(\mathbb{R}^m)$ 中成立. 为此需要下面称为 Friedrichs 引理的结果.

\begin{Lemma}{Friedrichs}
\label{lem:friedrichs-commutator-approximation}
对 $\epsilon>0$, 存在从 $L^2(\mathbb{R}^m)$ 到 $H^1(\mathbb{R}^m)$ 的映射 $Q_\epsilon$, 使以下性质成立:\begin{enumerate}\item 对每个 $\bm{f}\in L^2(\mathbb{R}^m)$, 当 $\epsilon\to0$ 时, $Q_\epsilon\bm{f}\to\bm{f}$ 在 $L^2(\mathbb{R}^m)$ 中成立.\item 算子 $\bm{A}Q_\epsilon-Q_\epsilon\bm{A}$ 可延拓为 $L^2(\mathbb{R}^m)$ 中的有界算子. 此外, 对任意固定的 $\bm{u}\in\mathcal{D}(\widetilde{\bm{A}})$, 有 $\bm{A}Q_\epsilon\bm{u}-Q_\epsilon\widetilde{\bm{A}}\bm{u}\to0$ 在 $L^2(\mathbb{R}^m)$ 中成立.\end{enumerate}
\end{Lemma}

有了这个引理, 证明立即完成. 只需置 $\bm{u}_\epsilon=Q_\epsilon\bm{u}$ 和 $\bm{f}_\epsilon=\bm{A}\bm{u}_\epsilon$. 则
\begin{equation}\label{eq:symmetric-hyperbolic-friedrichs-graph-approximation}
\bm{f}_\epsilon=Q_\epsilon\widetilde{\bm{A}}\bm{u}+(\bm{A}Q_\epsilon-Q_\epsilon\widetilde{\bm{A}})\bm{u}.
\end{equation}
当 $\epsilon\to0$ 时, 右端第一项收敛到 $\widetilde{\bm{A}}\bm{u}=\bm{f}$, 第二项收敛到零.

\begin{Proof}
\PDEIndexMark{idx:friedrichs-commutator-lemma-full-proof}因此只需再证明 Friedrichs 引理. 设 $k(\bm{x})$ 为支集位于球 $|\bm{x}|\leq1$ 内且积分为一的非负试验函数. 定义 $k_\epsilon(\bm{x})=\epsilon^{-m}k(\bm{x}/\epsilon)$, 并将 $Q_\epsilon$ 定义为与 $k_\epsilon$ 作卷积. 则对每个 $\epsilon>0$, $Q_\epsilon$ 将 $L^2(\mathbb{R}^m)$ 映入 $H^1(\mathbb{R}^m)$ (实际上映入 $C^\infty(\mathbb{R}^m)$); 此外, Lemma~\ref{lem:integral-kernel-row-column-boundedness} 蕴含, $Q_\epsilon$ 作为从 $L^2(\mathbb{R}^m)$ 到自身的算子, 范数不超过 1. 容易对 $L^2(\mathbb{R}^m)$ 的某个稠密子集证明 $Q_\epsilon\bm{u}\to\bm{u}$ (例如对 $\bm{u}\in\mathcal{D}(\mathbb{R}^m)$); 结合 $Q_\epsilon$ 的一致有界性, 可知 $Q_\epsilon$ 强收敛到恒等算子, 即引理的性质 1 成立.\par 为验证性质 2, 计算得到
\begin{equation}\label{eq:friedrichs-mollifier-commutator-integral-identity}
\begin{aligned}
\bm{A}Q_\epsilon\bm{u}-Q_\epsilon\widetilde{\bm{A}}\bm{u}&=\int_{\mathbb{R}^m}k_\epsilon(\bm{x}-\bm{y})(\bm{A}^i(\bm{x})-\bm{A}^i(\bm{y}))\dfrac{\partial\bm{u}}{\partial y_i}(\bm{y})\,\mathrm{d}\bm{y}\\
&=\int_{\mathbb{R}^m}k_\epsilon(\bm{x}-\bm{y})\dfrac{\partial\bm{A}^i}{\partial y_i}(\bm{y})\bm{u}(\bm{y})\,\mathrm{d}\bm{y}\\
&\quad+\int_{\mathbb{R}^m}\dfrac{\partial}{\partial x_i}(k_\epsilon(\bm{x}-\bm{y}))(\bm{A}^i(\bm{x})-\bm{A}^i(\bm{y}))\bm{u}(\bm{y})\,\mathrm{d}\bm{y}.
\end{aligned}
\end{equation}
当 $\epsilon\to0$ 时, 对足够光滑的 $\bm{u}$ (例如 $\bm{u}\in\mathcal{D}(\mathbb{R}^m)$), 右端第一个表达式显然趋于零. 因而只需证明, 当 $\epsilon\to0$ 时, $\bm{A}Q_\epsilon-Q_\epsilon\widetilde{\bm{A}}$ 作为从 $L^2(\mathbb{R}^m)$ 到自身的算子一致有界. 为此, 使用 \eqref{eq:friedrichs-mollifier-commutator-integral-identity} 右端的第二个表达式. 由 Lemma~\ref{lem:integral-kernel-row-column-boundedness} 的证明可知, 该表达式中第一项的算子范数不超过
\begin{equation}\label{eq:friedrichs-commutator-coefficient-derivative-norm-bound}
\sup_{\bm{y}\in\mathbb{R}^m}\left|\dfrac{\partial\bm{A}^i}{\partial y_i}(\bm{y})\right|\times\int_{\mathbb{R}^m}|k_\epsilon(\bm{x})|\,\mathrm{d}\bm{x},
\end{equation}
在这个表达式中, 第一个因子不依赖于 $\epsilon$, 第二个因子为 1. 再由 Lemma~\ref{lem:integral-kernel-row-column-boundedness}, 第二项的算子范数可用
\begin{equation}\label{eq:friedrichs-commutator-kernel-derivative-norm-bound}
\sup_{\bm{x}\in\mathbb{R}^m}\int_{\mathbb{R}^m}\left|\dfrac{\partial}{\partial x_i}k_\epsilon(\bm{x}-\bm{y})\right|\left|\bm{A}^i(\bm{x})-\bm{A}^i(\bm{y})\right|\,\mathrm{d}\bm{y}.
\end{equation}
估计. 由于 $k_\epsilon$ 的支集包含于 $|\bm{x}|\leq\epsilon$, 可以用某常数乘以 $\epsilon$ 估计 $|\bm{A}^i(\bm{x})-\bm{A}^i(\bm{y})|$. 另一方面, $\partial k_\epsilon/\partial x_i$ 的 $L^1$ 范数为某常数乘以 $1/\epsilon$. 因而 \eqref{eq:friedrichs-commutator-kernel-derivative-norm-bound} 有不依赖于 $\epsilon$ 的常数界.
\end{Proof}

\subsubsection{波动方程}\label{subsec:semigroup-wave-equation}

考虑二阶方程
\begin{equation}\label{eq:semigroup-wave-variable-coefficient-equation}
u_{tt}=\dfrac{\partial}{\partial x_i}\left(a_{ij}(\bm{x})\dfrac{\partial u}{\partial x_j}\right),\ \bm{x}\in\Omega,\ t>0,
\end{equation}
及边界条件
\begin{equation}\label{eq:semigroup-wave-dirichlet-boundary}
u(\bm{x},t)=0,\ \bm{x}\in\partial\Omega,
\end{equation}
和初始条件
\begin{equation}\label{eq:semigroup-wave-displacement-velocity-initial-data}
u(\bm{x},0)=u_0(\bm{x}),\ u_t(\bm{x},0)=u_1(\bm{x}).
\end{equation}
这里假设 $\Omega$ 为 $\mathbb{R}^m$ 中具有光滑边界的有界区域, 矩阵 $a_{ij}$ 对称, 严格正定且为 $C^1(\overline\Omega)$ 类.

用 $A$ 表示 \eqref{eq:semigroup-wave-variable-coefficient-equation} 右端给出的算子. 将 $A$ 视为 $L^2(\Omega)$ 中定义域为 $H^2(\Omega)\cap H_0^1(\Omega)$ 的算子. 为应用半群理论, 必须先将二阶方程 \eqref{eq:semigroup-wave-variable-coefficient-equation} 化为一阶系统. 为此置 $v=u_t$, 并将 \eqref{eq:semigroup-wave-variable-coefficient-equation} 改写为
\begin{equation}\label{eq:semigroup-wave-first-order-operator-matrix}
\begin{pmatrix}\dot u\\\dot v\end{pmatrix}=\mathcal{A}\begin{pmatrix}u\\v\end{pmatrix}:=\begin{pmatrix}0&I\\A&0\end{pmatrix}\begin{pmatrix}u\\v\end{pmatrix}.
\end{equation}
将 $\mathcal{A}$ 视为空间 $X=H_0^1(\Omega)\times L^2(\Omega)$ 中定义域为 $\mathcal{D}(\mathcal{A})=(H_0^1(\Omega)\cap H^2(\Omega))\times H_0^1(\Omega)$ 的算子. 断言

\begin{Theorem}
\label{thm:semigroup-wave-generation}
算子 $\mathcal{A}$ 生成 $C_0$ 半群.
\end{Theorem}

\begin{Proof}
再次使用 Lumer-Phillips 定理. 为此, 考虑 $X$ 上如下与通常内积等价的内积:
\begin{equation}\label{eq:semigroup-wave-energy-inner-product}
\left[\begin{pmatrix}u\\v\end{pmatrix},\begin{pmatrix}f\\g\end{pmatrix}\right]=\int_\Omega a_{ij}(\bm{x})\dfrac{\partial u}{\partial x_i}\dfrac{\partial f}{\partial x_j}+vg\,\mathrm{d}\bm{x}.
\end{equation}
容易看出, 对 $(u,v)\in\mathcal{D}(\mathcal{A})$,
\begin{equation}\label{eq:semigroup-wave-generator-energy-skew-pairing}
\left[\begin{pmatrix}u\\v\end{pmatrix},\mathcal{A}\begin{pmatrix}u\\v\end{pmatrix}\right]=0
\end{equation}
此外, 显然 $\mathcal{D}(\mathcal{A})$ 稠密. 断言每个正的 $\lambda$ 都属于 $\mathcal{A}$ 的预解集. 为此, 需要考虑问题
\begin{equation}\label{eq:semigroup-wave-resolvent-coupled-equations}
v-\lambda u=f,\ Au-\lambda v=g,\ f\in H_0^1(\Omega),\ g\in L^2(\Omega).
\end{equation}
这两个方程可以合并为一个方程
\begin{equation}\label{eq:semigroup-wave-resolvent-reduced-elliptic-equation}
Au-\lambda^2u=g+\lambda f.
\end{equation}
该方程的可解性由 Chapter~\ref{chap:linear-elliptic-equations} 中关于椭圆边值问题的结果得到.
\end{Proof}

\subsubsection{Schrödinger 方程}\label{subsec:semigroup-schrodinger-equation}

Schrödinger 方程是量子力学的基本方程. 它写为
\begin{equation}\label{eq:semigroup-schrodinger-real-potential-equation}
u_t=i(\Delta u-V(\bm{x})u),\ \bm{x}\in\mathbb{R}^m.
\end{equation}
这里 $i$ 为虚数单位, 即 $u$ 为复值函数. 然而, 函数 $V(\bm{x})$ 为实值函数. 在形式意义下, \eqref{eq:semigroup-schrodinger-real-potential-equation} 右端的算子斜对称; 若正确定义其定义域, 预期它为斜自伴算子. 在这种情形下, 立即可知该算子生成 $C_0$ 半群.\par 数学量子力学中相当大的一部分工作是证明这一算子的斜自伴性. 显然, 其难度取决于允许 $V$ 的行为有多差. $V$ 有界的情形基本上是平凡的, 但物理意义不大. 更有意义的是具有点奇性的势. 这里证明以下结果.

\begin{Theorem}
\label{thm:schrodinger-ltwo-potential-skew-adjointness}
假设 $m\leq3$ 且 $V\in L^2(\mathbb{R}^m)$. 则由 $Au=i(\Delta u-V(\bm{x})u)$ 给出且定义域为 $\mathcal{D}(A)=H^2(\mathbb{R}^m)$ 的算子 $A$ 为斜自伴算子.
\end{Theorem}

\begin{Proof}
注意到
\begin{equation}\label{eq:schrodinger-ltwo-potential-multiplication-bound}
\|Vu\|_2\leq\|V\|_2\|u\|_\infty\leq C\|V\|_2\|u\|_{m/2+\delta,2}\leq C\|V\|_2\|u\|_{2,2}.
\end{equation}
这里 $\delta$ 为任意足够小的正数. 这说明 $A$ 确实在 $H^2(\mathbb{R}^m)$ 上有定义. 此外, 利用 Fourier 变换作直接计算, 可将 \eqref{eq:schrodinger-ltwo-potential-multiplication-bound} 改进如下: 对每个 $\epsilon>0$, 存在常数 $C(\epsilon)$, 使
\begin{equation}\label{eq:schrodinger-potential-small-relative-interpolation-bound}
\|u\|_{m/2+\delta,2}\leq\epsilon\|u\|_{2,2}+C(\epsilon)\|u\|_2.
\end{equation}
利用这一点可证明 $A$ 为闭算子; 见 Problem~\ref{ex:schrodinger-potential-operator-closedness}. 显然 $A$ 为斜对称算子. 为证明它斜自伴, 需要证明 $A\pm I$ 为满射. 考虑 $t\in[0,1]$ 时的算子族 $i(\Delta-tV)\pm I$, 并利用 Fredholm 指标的同伦不变性, 即可得到这一点.
\end{Proof}

显然, 更一般地, 可以考虑为某个 $L^2$ 函数与某个有界函数之和的 $V$. 例如, 这包括 $\mathbb{R}^3$ 中的 Coulomb 势 ($V=-1/r$).

\Needspace{5\baselineskip}\textbf{Problems}

\begin{enumerate}[label=\thesection.\arabic*.,ref=\thesection.\arabic*,start=15,series=book-problems]

\item\label{ex:semigroup-transport-and-wave-existence-statements} 对带有初始条件 $u(x,y,0)=u_0(x,y)$ 的方程 $u_t=a(x,y)u_x+b(x,y)u_y+c(x,y)u+f(x,y,t)$, $(x,y)\in\mathbb{R}^2$, $t>0$, 陈述一个存在性定理. 类似地, 对波动方程陈述一个存在性定理.

\item\label{ex:strict-hyperbolic-system-symmetric-transformation} 考虑问题 $\bm{u}_t=\bm{A}(x)\bm{u}_x$, 其中矩阵 $\bm{A}$ 具有互异的实特征值. 证明这一问题总能化为具有对称矩阵的问题, 从而可应用 Section~\ref{subsec:semigroup-symmetric-hyperbolic-systems} 的结果. 在多个空间维度下会出现什么困难?

\item\label{ex:second-order-drift-operator-hthree-domain-generation} 在 $L^2(\mathbb{R})$ 中, 令 $Au=u_{xx}+a(x)u_x$, 其定义域为 $\mathcal{D}(A)=H^3(\mathbb{R})$. 假设 $a\in C_b^1(\mathbb{R})$. 证明 $A$ 生成 $C_0$ 半群.

\item\label{ex:wave-neumann-source-unresolved-refwhyp} 以类似于 Section \texttt{refWHYP} 的方式讨论带 Neumann 条件的波动方程.

\item\label{ex:schrodinger-potential-operator-closedness} 利用 \eqref{eq:schrodinger-potential-small-relative-interpolation-bound} 证明 Theorem~\ref{thm:schrodinger-ltwo-potential-skew-adjointness} 中的算子 $A$ 为闭算子.

\end{enumerate}
